\documentclass[cs_edp_2026.1.tex]{subfiles} % Aponta para o pai
\begin{document}
	
	% CONFIGURAÇÃO PARA COMPILAÇÃO INDIVIDUAL
	\ifSubfilesClassLoaded{
		\def\listid{L4}
		\setcounter{numquestao}{74}
		\section*{Equação da Onda} % Adiciona o título apenas na compilação isolada
	}{}
	
	%%%% QUESTÃO 75 %%%%
	\questao{(De volta a equação de transporte)
		\begin{enumerate}[label=(\alph*)]
			\item Suponha que se $v = v(x, t)$ satisfaz a equação de transporte homogênea $v_t + v_x = 0$ em $\mathbb{R} \times [0, \infty)$; $v(x, 0) = g(x); x \in \mathbb{R}$; então $v = g(x - t)$;
			\item Se $w = w(x, t)$ satisfaz a equação não homogênea de transporte $w_t - w_x = f(x, t)$ em $\mathbb{R} \times [0, \infty)$; $w(x, 0) = 0; x \in \mathbb{R}$; então 
			$$ w(x, t) = \int_0^t f(x + t - s, s) ds; $$
			(aqui tem o princípio de Duhamel)
			\item Se $f(x, t) = g(x - t)$ em (b), então 
			$$ w(x, t) = \int_0^t g(x + t - 2s) ds = \frac{1}{2} \int_{x-t}^{x+t} g(s) ds = \frac{1}{2} [G(x + t) - G(x - t)]; $$
			onde $G$ é uma primitiva de $g$.
			\item Considere $u$ uma solução de classe $C^2$ da equação da onda homogênea: $u_{tt} - u_{xx} = 0$ em $\mathbb{R} \times [0, \infty)$; $u(x, 0) = g(x); u_t(x, 0) = h(x); x \in \mathbb{R}$. Mostre que $v(x, t) = u_t - u_x$ satisfaz a equação de transporte do item (a), com $v(x, 0) = g'(x) + h(x)$, e conclua que $v(x, t) = g'(x - t) + h(x - t)$;
			\item Veja então que $u$ satisfaz a equação não homogênea $u_t - u_x = g'(x - t) + h(x - t)$ em $\mathbb{R} \times [0, \infty)$; $u(x, 0) = g(x)$;
			\item Conclua que 
			$$ u(x, t) = \frac{1}{2} \int_{x-t}^{x+t} h(s) ds + \frac{1}{2} [g(x + t) + g(x - t)]. $$
	\end{enumerate}}
	
	\analise{
		\textbf{Análise do Enunciado:} Esta questão foi estruturada como um roteiro de construção analítica. O objetivo não é apenas encontrar a solução final, mas derivar a fórmula de d'Alembert para a equação da onda através de etapas intermediárias. A lógica consiste em utilizar a teoria de equações de transporte (de 1\f ordem) para resolver a equação da onda (de 2\m ordem). Assim, cada item, de (a) a (e), atua como um lema necessário para a conclusão final no item (f), estabelecendo a ponte entre a propagação de ondas unidirecionais e a solução geral da equação da onda. As hipóteses centrais são a regularidade $u \in C^2$ no domínio $\mathbb{R} \times [0, \infty)$ e as condições iniciais clássicas $u(x,0) = g(x)$ e $u_t(x,0) = h(x)$.
		
		\par\textbf{Fundamentação Teórica:} Utilizamos o Método das Características (Evans, Cap. 2.2), onde a solução de $u_t + c u_x = 0$ é constante ao longo das retas $x - ct = \text{constante}$. Para a parte não homogênea, aplicamos o Princípio de Duhamel (Evans, Cap. 2.3), que trata a fonte externa como uma sucessão de impulsos que geram novas condições iniciais ao longo do tempo. A base para a redução da equação da onda é a fatoração do operador $\partial_{tt} - \partial_{xx}$ no produto $(\partial_t - \partial_x)(\partial_t + \partial_x)$.
		
		\par\textbf{Estratégias:} Começamos resolvendo a equação de transporte homogênea e a não homogênea via características. Em seguida, definimos a variável auxiliar $v = u_t - u_x$ para reduzir a ordem da equação da onda. Por fim, integramos a solução via Princípio da Superposição, somando a parte homogênea e a particular.
	}
	
	\solucao{
		\textbf{(a)} Começamos com a equação $v_t + v_x = 0$. Definimos a curva característica por $\tfrac{dx}{dt} = 1$. Integrando, temos $x(t) = t + x_0$, ou $x_0 = x - t$. Ao longo dessa curva, a variação de $v$ é:
		$$ \frac{d}{dt} v(x(t), t) = \frac{\partial v}{\partial x} \frac{dx}{dt} + \frac{\partial v}{\partial t} = v_x(1) + v_t = v_t + v_x $$
		Como a equação nos diz que $v_t + v_x = 0$, temos $\frac{d}{dt} v(x(t), t) = 0$, o que significa que $v$ é constante ao longo da característica. Assim, $v(x, t) = v(x_0, 0)$. Substituindo $x_0$ e a condição inicial $v(x,0) = g(x)$, chegamos a:
		$$ \boxed{v(x, t) = g(x - t) \quad (\ast)} $$
		
		\textbf{(b)} Agora analisamos $w_t - w_x = f(x, t)$ com $w(x,0) = 0$. As características são $\tfrac{dx}{ds} = -1$, logo $x(s) = x + t - s$. Calculando a derivada de $w$ ao longo dessa curva:
		$$ \frac{d}{ds} w(x + t - s, s) = -w_x + w_s = f(x + t - s, s) $$
		Integrando de $s=0$ até $s=t$:
		$$ \int_0^t \frac{d}{ds} w(x + t - s, s) ds = \int_0^t f(x + t - s, s) ds $$
		$$ w(x, t) - w(x+t, 0) = \int_0^t f(x + t - s, s) ds $$
		Como $w(x,0) = 0$, a solução é:
		$$ \boxed{w(x, t) = \int_0^t f(x + t - s, s) ds \quad (\ast\ast)} $$
		
		\textbf{(c)} Usamos o resultado $(\ast\ast)$ com a fonte $f(x, t) = g(x - t)$:
		$$ w(x, t) = \int_0^t g((x + t - s) - s) ds = \int_0^t g(x + t - 2s) ds $$
		Fazemos a mudança de variável $\sigma = x + t - 2s$, com $d\sigma = -2 ds$. Os limites mudam para $s=0 \to \sigma=x+t$ e $s=t \to \sigma=x-t$. A integral fica:
		$$ w(x, t) = \int_{x+t}^{x-t} g(\sigma) \left( -\tfrac{1}{2} \right) d\sigma = \tfrac{1}{2} \int_{x-t}^{x+t} g(\sigma) d\sigma = \tfrac{1}{2} [G(x + t) - G(x - t)] $$
		
		\textbf{(d)} Seja $v = u_t - u_x$. Calculamos as derivadas parciais:
		$$ v_t = u_{tt} - u_{tx} \quad \text{e} \quad v_x = u_{tx} - u_{xx} $$
		Somando as duas, temos $v_t + v_x = u_{tt} - u_{xx}$. Como $u$ resolve a equação da onda homogênea, $u_{tt} - u_{xx} = 0$, logo $v_t + v_x = 0$. Para a condição inicial:
		$$ v(x, 0) = u_t(x, 0) - u_x(x, 0) = h(x) - g'(x) $$
		\textit{Nota: O PDF indica $g'(x)+h(x)$, mas a dedução rigorosa dá $h(x)-g'(x)$. Usaremos o sinal correto para a consistência de (f).}
		
		\vspace{0.2cm}
		Observamos que a equação $v_t + v_x = 0$ com dado inicial $v(x,0) = h(x) - g'(x)$ é idêntica à estrutura resolvida no item (a). Portanto, aplicando a solução $(\ast)$, temos:
		$$ v(x, t) = h(x - t) - g'(x - t) $$
		
		\textbf{(e)} Pela definição de $v$ no item (d), temos:
		$$ u_t - u_x = v(x, t) = h(x - t) - g'(x - t) $$
		com $u(x, 0) = g(x)$. Esta é a equação de transporte não homogênea do item (b), cuja solução geral é dada por $(\ast\ast)$.
		
		\textbf{(f)} Para achar $u$, usamos o Princípio da Superposição, somando a solução homogênea $u_h$ e a particular $u_p$ ($u = u_h + u_p$).
		Para $u_h$, resolvemos $u_t - u_x = 0$ com $u(x,0) = g(x)$. Pelas características $\tfrac{dx}{dt} = -1$, temos $u_h(x, t) = g(x + t)$.
		Para $u_p$, resolvemos $u_t - u_x = f$ com $u(x,0) = 0$. Usando o resultado $(\ast\ast)$ com $f(x, s) = h(x - s) - g'(x - s)$:
		$$ u_p = \int_0^t [h(x + t - 2s) - g'(x + t - 2s)] ds $$
		Fazendo $\sigma = x + t - 2s$ ($ds = -\tfrac{1}{2} d\sigma$), temos:
		$$ u_p = \tfrac{1}{2} \int_{x-t}^{x+t} h(\sigma) d\sigma - \tfrac{1}{2} \int_{x-t}^{x+t} g'(\sigma) d\sigma = \tfrac{1}{2} \int_{x-t}^{x+t} h(\sigma) d\sigma - \tfrac{1}{2} [g(x + t) - g(x - t)] $$
		Somando $u_h$ e $u_p$:
		$$ u(x, t) = g(x + t) + \tfrac{1}{2} \int_{x-t}^{x+t} h(s) ds - \tfrac{1}{2} g(x + t) + \tfrac{1}{2} g(x - t) $$
		Simplificando, chegamos à fórmula de d'Alembert:
		$$ \boxed{u(x, t) = \tfrac{1}{2} [g(x + t) + g(x - t)] + \tfrac{1}{2} \int_{x-t}^{x+t} h(s) ds \quad (\ast\ast\ast)} $$
	}
	
	\sintese{
		A resolução desta questão demonstra que a equação da onda é a superposição de dois fluxos de transporte em direções opostas. A \textbf{Construção Analítica} do raciocínio pode ser resumida no seguinte fluxo:
		
		\begin{center}
			$\text{Item (a)} \longrightarrow \text{Item (d)} \longrightarrow \text{Item (e)} \longrightarrow \text{Item (f)}$ \\
			$\text{Item (b)} \longrightarrow \text{Item (f)}$
		\end{center}
		
		O resultado do Item (a) fornece a técnica de solução para a direita, usada no Item (d) para descobrir a variável auxiliar $v$. O resultado do Item (b) fornece a solução para o transporte não homogêneo, usada no Item (f) para reconstruir a solução completa de d'Alembert.
	}
	
\end{document}